\section{编译与运行结果}

\subsection{运行环境}

\begin{table}[htbp]
\centering
\caption{运行与检查环境}\label{tab:env}
\zihao{-5}
\renewcommand{\arraystretch}{1.15}
\begin{tabularx}{\textwidth}{@{}L{3.0cm}L{5.0cm}Y@{}}
\toprule
项目 & 版本 / 型号 & 说明 \\
\midrule
操作系统 & Windows 11（win32） & 命令行使用 Git Bash \\
Python & 3.13.14 & 仅标准库；无运行时第三方依赖 \\
代码检查 & pylint 4.1.3 & PEP 8，行宽 79 \\
单元测试 & pytest 9.2.0 & 73 个用例 \\
UML 渲染 & PlantUML 1.2024.8 + OpenJDK 17 & 输出 190 dpi PNG \\
报告排版 & TeX Live 2025（XeLaTeX + ctex） & —— \\
\bottomrule
\end{tabularx}
\end{table}

\subsection{编译检查}

Python 是解释型语言，没有独立的编译步骤；为满足"使用编译器或 IDE 编译"的要求，用标准库 \code{compileall} 对全部源码做字节码编译，它会同时暴露语法错误与模块导入错误：

\begin{lstlisting}
$ python -m compileall -q campustrade run.py tests
$ echo $?
0
\end{lstlisting}

没有任何输出、退出码为 0，说明 22 个源文件与 8 个测试文件都能通过语法与导入检查。运行时验证由 \code{python run.py --demo}（端到端）与交互式 \code{python run.py --seed} 两条路径覆盖，下面分别给出真实输出。

\subsection{单元测试}

\begin{lstlisting}
$ python -m pytest -q
\end{lstlisting}

\lstinputlisting{listings/pytest.tex}

测试按模块分文件，覆盖关系见\tabref{tab:tests}。这里刻意没有追求行覆盖率数字，而是每写一条规则就配一个"应该被拒绝"的用例：权限越权、状态不允许、重复提交、时间未到，这些分支才是真实缺陷的藏身处。

\begin{table}[htbp]
\centering
\caption{测试文件与覆盖的需求}\label{tab:tests}
\zihao{-5}
\renewcommand{\arraystretch}{1.15}
\begin{tabularx}{\textwidth}{@{}L{3.4cm}L{1.4cm}Y@{}}
\toprule
测试文件 & 用例数 & 覆盖内容 \\
\midrule
\code{test\_auth.py} & 12 & 注册、登录、重复学号、弱口令、会话过期、停用账号、改密码（REQ-004/012/016） \\
\code{test\_item.py} & 15 & 草稿、公开查询过滤与排序、必填校验、下架与恢复、分类唯一性与引用保护（REQ-002/005/009/014） \\
\code{test\_audit.py} & 10 & 通过/驳回、待审内容原子替换、驳回丢弃、版本冲突、审核历史与可见性（REQ-003/017/019） \\
\code{test\_trade.py} & 13 & 幂等重发、唯一有效预约、取消、确认、完成、超时自动完成、异议暂停、并发竞争（REQ-001/006/007/008/015） \\
\code{test\_moderation.py} & 8 & 举报去重、处置下架/驳回、重复处理、已预约商品不能下架（REQ-010/014） \\
\code{test\_storage.py} & 10 & JSON 快照往返（Decimal/枚举/时间）、时间范围统计、日志权限、损坏文件处理（REQ-013/018） \\
\code{test\_cli.py} & 5 & 登录、搜索、预约、非法选项、管理员菜单的界面冒烟 \\
\midrule
合计 & 73 & —— \\
\bottomrule
\end{tabularx}
\end{table}

\subsection{端到端演示}

\code{python run.py --demo} 在内存库上按需求编号跑 11 节场景，每节既有"应当成功"的路径，也有"应当被拒绝"的反例，最后打印是否出现未捕获异常。完整输出如下（\code{√} 表示按预期成功，\code{×} 表示按预期被拒绝）：

\lstinputlisting{listings/demo.tex}

几个值得单独说明的结果：

\begin{itemize}[leftmargin=1.6em,itemsep=2pt,topsep=3pt]
\item 第 4 节：同一 \code{requestId} 重发返回预约号 1，预约记录仍只有 1 条，"重发不算新请求"落地；另一买家再预约被拒，唯一性成立。
\item 第 7 节：修改已发布商品后前台仍显示《罗技 M590 静音鼠标》89 元，期间买家仍能预约（预约号 3），审核通过后价格变为 79 元且待审记录被清除；驳回的那一件仍显示原价 45 元。
\item 第 9 节：4 个买家经屏障同步后同时预约，只有 1 个成功，其余 3 个收到 \code{CONFLICT}。商品最终只有 1 条有效预约。
\item 第 10 节：把确认时间回拨 8 天后触发自动完成，商品变已完成，卖家与买家都收到通知。
\end{itemize}

\subsection{交互式运行}

\code{python run.py --seed} 进入菜单界面，下面是真实会话。为节省篇幅，省去了每条命令后重复打印的菜单与输入提示，只在 \code{→} 后保留该步的执行结果，其余原样保留：管理员审核通过一件待审商品，买家搜索并预约，卖家确认交易，买家确认完成，最后卖家查看统计。

\lstinputlisting{listings/cli.tex}

会话中可以看到权限是按身份变化的：管理员登录后菜单多出 21--30 项（审核、举报、统计、日志、分类），普通用户只有 1--15 项；未登录时只能浏览商品。最后一行"发布 4 件，预约 1 次，完成 1 笔，交易额 45.00 元"与前面几步的操作完全对得上，说明统计口径没有把别人的数据算进来。
